\documentclass[10pt,a4paper]{amsart}
\usepackage[margin=19mm]{geometry}
\usepackage{amsmath,amssymb,mathtools}
\usepackage[scr=rsfs,cal=euler]{mathalfa}
\usepackage{xfrac}
\usepackage[unicode,hidelinks,bookmarks=false]{hyperref}
\newcommand{\ie}{i.e.\ }
\newcommand{\cf}{cf.\ }
\newcommand{\resp}{respectively\ }
\newcommand{\Subsection}{\subsection}
\providecommand{\nobreakdash}{\nobreak\hbox{-}}
\setlength{\emergencystretch}{2em}
\multlinegap0pt
\usepackage{bm}
\usepackage{bbm}
\usepackage{dsfont}
\usepackage{xspace}
\usepackage{cancel}
\usepackage{mathtools}
\usepackage{amsthm}
\makeatletter
\@ifundefined{theorem}{\newtheorem{theorem}{Theorem}}{}
\@ifundefined{proposition}{\newtheorem{proposition}[theorem]{Proposition}}{}
\makeatother
\title[The Quantum Adiabatic Theorem for Non-Hermitian Dynamics]{The Quantum Adiabatic Theorem for Non-Hermitian Dynamics}
\author{G. M. V. S. Aponso, B. P. W. Fernando, K. K. W. A. S. Kumara, R. P. K. De Silva}
\date{Source: arXiv:2607.19737v1 (CC BY 4.0)}
\begin{document}
\maketitle

\noindent\emph{Source and licence.} The Quantum Adiabatic Theorem for Non-Hermitian Dynamics. Version arXiv:2607.19737v1 (CC BY 4.0), distributed under the Creative Commons Attribution 4.0 International licence, as recorded in the arXiv licence metadata for this exact version and shown on the abstract page https://arxiv.org/abs/2607.19737v1. Full rights notice: licenses/arxiv-2607.19737-CC-BY-4.0.txt.\par\smallskip

\noindent\textit{Editorial scope.} Counted unit: the Proposition whose source label is \texttt{prop:ProjectionComparison} in the article (source0.tex lines 940--965, source label \texttt{prop:ProjectionComparison}), delivered together with the article's own complete original proof of it (source0.tex lines 967--1068, one \texttt{proof} environment of 102 lines). No mathematics is written, repaired or re-derived: statement and proof are the article's own text line for line. Required context retained uncounted: none. Every definition, display and label of those lines is retained unchanged. Omitted from this package and not redistributed: 1--939, 966--966, 1069--1837. Nothing inside a retained excerpt is omitted or altered: every retained body is delivered exactly as the article's lines are written, so no omission receipt of any kind accompanies this package. Citations: the retained excerpts carry no citation command at all, so no bibliography entry of the article is transcribed and no reference is redistributed. Reference adaptation: none was needed and none was made; every reference inside the delivered unit resolves inside it, so no unresolved reference or citation is produced. Printed slips kept literal: no printed word, symbol, hypothesis or number of the counted statement or of its proof was altered. Layout and class: the article's own class is \texttt{sn-jnl} with packages including graphicx, multirow, amsmath, amssymb, amsfonts, amsthm, mathrsfs, appendix, xcolor, textcomp, manyfoot, booktabs, algorithm, algorithmicx; the wrapper is the lane's pinned \texttt{amsart} editorial wrapper at 10pt on A4, which restates the theorem-like environments the retained text needs and adds the article's own macro definitions that the retained text needs (none), with extra wrapper packages (bm, bbm, dsfont, xspace, cancel, mathtools). The delivered numbering is the unit's own. Assets: no retained span contains a figure environment, a graphics-inclusion command, a PDF-page inclusion, a table or a TikZ picture; no image file is copied into this package, the package declares no asset dependency and no image file is redistributed with it. Licence: version arXiv:2607.19737v1 of this article is distributed under the Creative Commons Attribution 4.0 International licence, as recorded in the arXiv licence metadata for this exact version and shown on the abstract page https://arxiv.org/abs/2607.19737v1. A full rights notice is delivered at licenses/arxiv-2607.19737-CC-BY-4.0.txt. Characters: every retained excerpt is pure ASCII, so the delivered engine, XeLaTeX with its default Latin Modern text font, reports no missing character.
\begin{proposition}[Comparison of spectral projections]
\label{prop:ProjectionComparison}
Let \(\hat P_j(s)\) be the spectral projection of \(\hat H_T(s)\), and let,
\[
\Pi_j(s):=\hat\eta_T(s)^{-1}\bar P_j(s)\hat\eta_T(s)
\]
be the Dyson-pulled-back spectral projection associated with \(\bar H_T(s)\).
Let \(\Gamma_j(s)\) be a positively oriented contour enclosing only the eigenvalue
\(\lambda_j(s)\) of \(A_T(s)\), and define,
\[
M_j(s):=\sup_{z\in\Gamma_j(s)}\|(zI-A_T(s))^{-1}\|.
\]
Assume that,
\[
M_j(s)\|B_T(s)\|<1.
\]
Then,
\[
\|\Pi_j(s)-\hat P_j(s)\|
\le
\|\hat\eta_T(s)^{-1}\|\,\|\hat\eta_T(s)\|\,
\frac{|\Gamma_j(s)|}{2\pi}
\frac{M_j(s)^2\|B_T(s)\|}
{(1-M_j(s)\|B_T(s)\|)}.
\]
\end{proposition}

\begin{proof}
Let \(Q_j(s)\) denote the spectral projection of,
\[
A_T(s)=\hat\eta_T(s)\hat H_T(s)\hat\eta_T(s)^{-1}
\]
associated with the eigenvalue \(\lambda_j(s)\). Since \(A_T(s)\) is similar to
\(\hat H_T(s)\), the corresponding spectral projections satisfy,
\[
Q_j(s)=\hat\eta_T(s)\hat P_j(s)\hat\eta_T(s)^{-1}.
\]

Since \(M_j(s)\|B_T(s)\|<1\), the contour \(\Gamma_j(s)\) lies in the
resolvent set of \(A_T(s)+\tau B_T(s)\) for every \(\tau\in[0,1]\).
Consequently, no spectrum crosses the contour during the perturbation from
\(A_T(s)\) to \(\bar H_T(s)=A_T(s)+B_T(s)\). We denote by \(\bar P_j(s)\) the
spectral projection of \(\bar H_T(s)\) associated with the spectral set enclosed by
\(\Gamma_j(s)\).

The spectral projection of \(\bar H_T(s)=A_T(s)+B_T(s)\) is given by the formula,
\[
\bar P_j(s)=\frac{1}{2\pi i}\int_{\Gamma_j(s)}
(zI-\bar H_T(s))^{-1}\,dz,
\]
whereas,
\[
Q_j(s)=\frac{1}{2\pi i}\int_{\Gamma_j(s)}
(zI-A_T(s))^{-1}\,dz.
\]
Therefore,
\[
\bar P_j(s)-Q_j(s)
=
\frac{1}{2\pi i}\int_{\Gamma_j(s)}
\left[(zI-\bar H_T(s))^{-1}-(zI-A_T(s))^{-1}\right]dz.
\]
Using the resolvent identity,
\[
(zI-\bar H_T(s))^{-1}-(zI-A_T(s))^{-1}
=
(zI-\bar H_T(s))^{-1}B_T(s)(zI-A_T(s))^{-1},
\]
we obtain,
\[
\|\bar P_j(s)-Q_j(s)\|
\le
\frac{|\Gamma_j(s)|}{2\pi}
\sup_{z\in\Gamma_j(s)}
\|(zI-\bar H_T(s))^{-1}\|\,\|B_T(s)\|\,
\|(zI-A_T(s))^{-1}\|.
\]

Since,
\[
zI-\bar H_T(s)
=
zI-A_T(s)-B_T(s),
\]
we can write,
\[
zI-\bar H_T(s)
=
\left(I-B_T(s)(zI-A_T(s))^{-1}\right)(zI-A_T(s)).
\]
The assumption \(M_j(s)\|B_T(s)\|<1\) implies that,
\(
I-B_T(s)(zI-A_T(s))^{-1}
\)
is invertible for all \(z\in\Gamma_j(s)\). By the Neumann series estimate,
\[
\|(zI-\bar H_T(s))^{-1}\|
\le
\frac{M_j(s)}{(1-M_j(s)\|B_T(s)\|)}.
\]
Therefore,
\[
\|\bar P_j(s)-Q_j(s)\|
\le
\frac{|\Gamma_j(s)|}{2\pi}
\frac{M_j(s)^2\|B_T(s)\|}{(1-M_j(s)\|B_T(s)\|)}.
\]

Finally, since,
\[
\Pi_j(s)=\hat\eta_T(s)^{-1}\bar P_j(s)\hat\eta_T(s),
\qquad
\hat P_j(s)=\hat\eta_T(s)^{-1}Q_j(s)\hat\eta_T(s),
\]
we have,
\[
\Pi_j(s)-\hat P_j(s)
=
\hat\eta_T(s)^{-1}\big(\bar P_j(s)-Q_j(s)\big)\hat\eta_T(s).
\]
Taking norms gives,
\[
\|\Pi_j(s)-\hat P_j(s)\|
\le
\|\hat\eta_T(s)^{-1}\|\,\|\hat\eta_T(s)\|\,
\|\bar P_j(s)-Q_j(s)\|.
\]
This proves the estimate.
\end{proof}

\end{document}
